\documentclass[11pt, letterpaper]{article}

% ── Packages ──────────────────────────────────────────────────────────────────
\usepackage[margin=1in]{geometry}
\usepackage{amsmath, amssymb}
\usepackage{hyperref}
\usepackage{xcolor}
\usepackage{booktabs}
\usepackage{enumitem}
\usepackage{titlesec}
\usepackage{parskip}
\usepackage{fancyhdr}
\usepackage[T1]{fontenc}

% ── Colour palette ─────────────────────────────────────────────────────────────
\definecolor{trinityblue}{RGB}{10, 60, 130}
\definecolor{darkgray}{RGB}{60, 60, 60}
\definecolor{lightgray}{RGB}{245, 245, 245}
\definecolor{accentblue}{RGB}{30, 100, 200}

% ── Hyperref setup ─────────────────────────────────────────────────────────────
\hypersetup{
  colorlinks=true,
  linkcolor=trinityblue,
  urlcolor=accentblue,
  citecolor=trinityblue
}

% ── Section title formatting ───────────────────────────────────────────────────
\titleformat{\section}
{\large\bfseries\color{trinityblue}}
{}
{0em}
{}
[\vspace{-0.5ex}\rule{\linewidth}{0.8pt}\vspace{0.5ex}]

\titleformat{\subsection}
{\normalsize\bfseries\color{darkgray}}
{}
{0em}
{}

\titlespacing{\section}{0pt}{1.2em}{0.5em}
\titlespacing{\subsection}{0pt}{0.8em}{0.3em}

% ── Header / Footer ────────────────────────────────────────────────────────────
\pagestyle{fancy}
\fancyhf{}
\renewcommand{\headrulewidth}{0pt}
\fancyfoot[L]{\small\color{darkgray} CPSC 498 / ECE Senior Project}
\fancyfoot[R]{\small\color{darkgray}\thepage}

% ── List style ─────────────────────────────────────────────────────────────────
\setlist[itemize]{noitemsep, topsep=2pt, leftmargin=1.4em}
\setlist[enumerate]{noitemsep, topsep=2pt, leftmargin=1.4em}

% ==============================================================================
\begin{document}

% ── Title Block ───────────────────────────────────────────────────────────────
\begin{center}
  {\LARGE\bfseries\color{trinityblue} Edge-Compute Pacing \& Feedback Node}\\[0.5em]
  {\large\color{darkgray} CPSC 498 / ECE Senior Project --- Concept Summary}\\[0.8em]
  \begin{tabular}{rl}
    \textbf{Students:}         & Sean Balbale \& Zoe Davidow \\
    \textbf{Program:}          & B.S.\ ECE\,/\,B.S.\ CS, Trinity College (Class of 2027) \\
    \textbf{Proposed Advisor:} & Prof.\ Kenneth A.\ Kousen \\
    \textbf{Date:}             & April 28, 2026
  \end{tabular}
\end{center}

\vspace{0.3em}
\noindent\rule{\linewidth}{1.2pt}
\vspace{0.3em}

% ── Abstract ──────────────────────────────────────────────────────────────────
\section{Abstract}

This project proposes the development of an embedded edge-compute node designed to provide
real-time biomechanical and fatigue feedback to endurance athletes --- specifically rowers and
cyclists. The device interfaces with existing wireless telemetry (BLE\,/\,ANT+) to aggregate
live sensor data from power meters and heart-rate monitors, executes a physiologically grounded
fatigue-prediction algorithm ($W'$ balance) directly on a resource-constrained microcontroller,
and translates the result into an intuitive visual interface via a programmable LED array.
The core research contribution is architectural: demonstrating rigorous, low-latency numerical
computation on the edge --- without offloading to a smartphone or cloud --- under strict power
and memory budgets.

% ── Problem Statement ─────────────────────────────────────────────────────────
\section{Problem Statement}

Commercial cycling and rowing head units display raw telemetry (watts, cadence, split) but
perform no predictive computation. Athletes must self-interpret numbers while under physiological
stress, a cognitively expensive task that degrades pacing decisions.

Smartphones and cloud offload solve this but introduce \textit{connectivity dependency},
\textit{latency}, and \textit{battery vulnerability} --- unacceptable in competitive race
settings. The problem this project addresses is therefore:

\begin{center}
  \colorbox{lightgray}{\parbox{0.88\linewidth}{
      \vspace{0.4em}
      \textit{Can a real-time fatigue model be implemented with sub-100\,ms latency on a
        bare-metal or RTOS microcontroller, with sufficient numerical precision to provide
      actionable pacing feedback --- entirely offline?}
      \vspace{0.4em}
  }}
\end{center}

% ── The W-Prime Model ─────────────────────────────────────────────────────────
\section{Core Algorithm: The \texorpdfstring{$W'$}{W-Prime} Balance Model}

\subsection{Physiological Background}

The $W'$ (W-prime) model, formalised by Monod \& Scherrer (1965) and extended by Skiba et al.\
(2012), describes an athlete's finite anaerobic work capacity above their \textit{Critical Power}
($CP$) threshold.  Two athlete-specific parameters govern the model:

\begin{itemize}
  \item $CP$ --- the highest sustainable steady-state power output (watts), analogous to
    lactate threshold.
  \item $W'$ --- the total finite energy reservoir (joules) available for supra-threshold
    efforts.
\end{itemize}

\subsection{Continuous Depletion (Monod--Scherrer)}

When power output $P(t) > CP$, the reservoir depletes. The instantaneous balance is:

\begin{equation}
  W'_{\mathrm{bal}}(t)
  \;=\;
  W' \;-\; \int_{0}^{t} \bigl[P(\tau) - CP\bigr]^{+} \, d\tau
  \label{eq:depletion}
\end{equation}

where $[x]^{+} = \max(x,\,0)$ is the positive-part operator, ensuring only supra-threshold
power draws from the reservoir.

\subsection{Reconstitution Model (Skiba et al., 2012)}

When $P(t) < CP$, the reservoir recovers exponentially. Skiba's differential formulation
unifies both phases in a single ODE:

\begin{equation}
  \frac{dW'_{\mathrm{bal}}}{dt}
  \;=\;
  \begin{cases}
    -\!\bigl(P(t) - CP\bigr)
    & \text{if } P(t) \geq CP \quad (\text{depletion}) \\[6pt]
    \dfrac{W' - W'_{\mathrm{bal}}(t)}{\tau_{W'}}
    & \text{if } P(t) < CP \quad (\text{recovery})
  \end{cases}
  \label{eq:skiba}
\end{equation}

The recovery time constant $\tau_{W'}$ is itself power-dependent:

\begin{equation}
  \tau_{W'}(t)
  \;=\;
  546\,e^{-0.01\,[CP - P(t)]} \;+\; 316
  \quad [\text{seconds}]
  \label{eq:tau}
\end{equation}

This yields fast recovery near $CP$ and slow recovery at very low intensities, consistent
with observed exercise physiology.

\subsection{Discrete-Time Implementation (On-Device)}

The MCU samples sensors at fixed intervals $\Delta t$ (target: $\Delta t = 1\,\mathrm{s}$
from the ANT+ broadcast). We replace the continuous integral with a first-order Euler
update:

\begin{equation}
  W'_{\mathrm{bal}}[n]
  \;=\;
  \begin{cases}
    W'_{\mathrm{bal}}[n-1]
    \;-\;\bigl(P[n] - CP\bigr)\,\Delta t
    & P[n] \geq CP \\[8pt]
    W'_{\mathrm{bal}}[n-1]
    \;+\;\dfrac{W' - W'_{\mathrm{bal}}[n-1]}{\tau_{W'}[n]}\,\Delta t
    & P[n] < CP
  \end{cases}
  \label{eq:discrete}
\end{equation}

with boundary condition $W'_{\mathrm{bal}}[0] = W'$ (athlete fully recovered at start).
This requires only fixed-point arithmetic per tick --- well within the Cortex-M4F budget
at 64\,MHz.

\subsection{LED Feedback Mapping}

The scalar output $W'_{\mathrm{bal}}[n] \in [0,\,W']$ is normalised and mapped to an
$N$-pixel LED strip:

\begin{equation}
  k_{\mathrm{lit}}[n]
  \;=\;
  \left\lfloor N \cdot \frac{W'_{\mathrm{bal}}[n]}{W'} \right\rfloor,
  \qquad
  \mathrm{hue}[n] =
  \begin{cases}
    \text{green} & k_{\mathrm{lit}}/N > 0.50 \\
    \text{amber} & 0.20 < k_{\mathrm{lit}}/N \leq 0.50 \\
    \text{red}   & k_{\mathrm{lit}}/N \leq 0.20
  \end{cases}
  \label{eq:led}
\end{equation}

A haptic / buzzer interrupt fires when $W'_{\mathrm{bal}}[n] < 0.10\cdot W'$.

% ── Technical Architecture ────────────────────────────────────────────────────
\section{Technical Architecture}

\subsection{Hardware Layer}
\begin{itemize}
  \item \textbf{MCU:} Nordic nRF52840 --- native dual-mode BLE\,5 + ANT+ radio,
    ARM Cortex-M4F @ 64\,MHz, 256\,kB RAM, 1\,MB Flash.
  \item \textbf{Power:} 18650 Li-Ion cell + BMS; target runtime $\geq 6$\,h.
  \item \textbf{Output:} WS2812B NeoPixel strip (8--12 LEDs) + piezo buzzer.
  \item \textbf{Form factor:} Handlebar mount (cycling) or hull-mount (rowing shell).
\end{itemize}

\subsection{Firmware \& Software Layer}
\begin{itemize}
  \item \textbf{RTOS:} FreeRTOS --- sensor task, compute task, output task with
    defined priorities and circular-buffer IPC.
  \item \textbf{Protocols:} ANT+ Bicycle Power Profile (PWR) and Heart Rate Profile
    (HRM); BLE GATT power service as fallback.
  \item \textbf{Language:} C\,/\,C++; algorithm unit-tested against reference
    Python implementation before deployment.
  \item \textbf{Latency target:} sensor-to-LED update $\leq 100\,\mathrm{ms}$
    end-to-end.
\end{itemize}

% ── CS Research Framing ───────────────────────────────────────────────────────
\section{Computer Science Research Framing}

The hardware is a \emph{constraint}, not the contribution.  The open CS research questions are:

\begin{enumerate}
  \item \textbf{Algorithmic correctness under discretisation.}  How does Euler-step
    sampling frequency affect $W'_{\mathrm{bal}}$ accuracy relative to the
    continuous-time ground truth in Eq.\,\eqref{eq:depletion}?

  \item \textbf{Concurrent firmware architecture.}  What FreeRTOS scheduling policy
    minimises worst-case latency while preventing sensor-buffer overflow on a
    single-core MCU?

  \item \textbf{Extensibility toward ML.}  Can a lightweight model (e.g., tiny LSTM)
    trained on historical power traces improve personalisation of $\tau_{W'}$
    beyond the population-level constants in Eq.\,\eqref{eq:tau}?
\end{enumerate}

These questions sit at the intersection of \textit{computational methods},
\textit{concurrent systems design}, and \textit{AI applied to software} ---
directly aligned with Prof.\ Kousen's published research interests.

% ── Deliverables & Timeline ───────────────────────────────────────────────────
\section{Deliverables \& Proposed Timeline}

\begin{center}
  \begin{tabular}{lll}
    \toprule
    \textbf{Phase} & \textbf{Milestone} & \textbf{Target} \\
    \midrule
    1 & Hardware bring-up; BLE\,/\,ANT+ sensor pairing verified & Sep 2026 \\
    2 & $W'$ algorithm on MCU; validated against Python reference & Oct 2026 \\
    3 & RTOS integration; latency and scheduling benchmarks & Nov 2026 \\
    4 & Full feedback loop; field test with live power meter & Dec 2026 \\
    5 & Formal write-up, poster, and final public demo & Apr 2027 \\
    \bottomrule
  \end{tabular}
\end{center}

% ── What We're Asking For ─────────────────────────────────────────────────────
\section{What We Are Asking of an Advisor}

\begin{itemize}
  \item \textbf{Monthly check-ins} to review software architecture and algorithmic
    design decisions.
  \item \textbf{Feedback on written deliverables} --- proposal, interim report, and
    final thesis --- well-suited to Prof.\ Kousen's background as a published
    technical author.
  \item \textbf{Guidance on CS rigour:} framing the project's contribution to satisfy
    CPSC 498 standards, particularly on the algorithmic and systems-design components.
  \item \textbf{Optional:} advising on any ML-based extensions (personalised
    $\tau_{W'}$ estimation) pursued in Phase 5.
\end{itemize}

Deep embedded-hardware expertise is \textbf{not} required --- hardware selection is
settled.  The open problems are software architecture and algorithmic design.

\vspace{1em}
\noindent\rule{\linewidth}{0.5pt}
\vspace{0.2em}
{\small
  \noindent \textbf{Contact:}\;
  Sean Balbale ---
  \href{mailto:sean.balbale@trincoll.edu}{sean.balbale@trincoll.edu}
  $\;\big|\;$
  Zoe Davidow ---
  \href{mailto:zoe.davidow@trincoll.edu}{zoe.davidow@trincoll.edu}
}

\end{document}
